% Presentación oral del TP5 (problema 2). Mismo formato Beamer que el TP4.
\documentclass[aspectratio=169,11pt]{beamer}
\usetheme{Warsaw}
\useoutertheme{miniframes}
\setbeamertemplate{navigation symbols}{}
\setbeamertemplate{footline}[frame number]
\let\originalpercent\%
\usepackage[spanish,es-noquoting]{babel}
\AtBeginDocument{\let\%\originalpercent}
\usepackage{amsmath,amssymb,graphicx,booktabs,xcolor,tikz}
\usetikzlibrary{positioning,arrows.meta,shapes.multipart}

\definecolor{itbablue}{RGB}{0,91,145}
\definecolor{accent}{RGB}{211,92,32}
\definecolor{softblue}{RGB}{232,243,249}
\definecolor{evento}{HTML}{2A78D6}
\definecolor{paso}{HTML}{EB6834}
\setbeamercolor{structure}{fg=itbablue}
\setbeamercolor{alerted text}{fg=accent}
\setbeamercolor{block title}{fg=white,bg=itbablue}
\setbeamercolor{block body}{bg=softblue}
\setbeamerfont{frametitle}{size=\large}
\setbeamerfont{title}{size=\LARGE}
\setbeamertemplate{section page}{\vfill\centering{\usebeamercolor[fg]{structure}\usebeamerfont{title}\insertsectionhead\par}\vfill}
\graphicspath{{figuras/}}

\InputIfFileExists{valores.tex}{}{}
\InputIfFileExists{enlaces.tex}{}{}

\newcommand{\vect}[1]{\mathbf{#1}}
\newcommand{\dt}{\Delta t}
\newcommand{\tstar}{T^{*}}
% Valores que todavía no están (los define valores.tex cuando existen sus resultados).
\providecommand{\pendiente}{\textcolor{accent}{pendiente}}
\newcommand{\valor}[1]{\ifcsname #1\endcsname\csname #1\endcsname\else\pendiente\fi}
\newcommand{\enlace}[1]{\ifcsname enlace#1\endcsname\csname enlace#1\endcsname\else Enlace de YouTube pendiente\fi}
% Grupos de texto sin fondo ni borde, con títulos y espaciado comunes.
\newcommand{\ficha}[3][itbablue]{%
  \par
  \begingroup
    \begin{minipage}{\linewidth}
    {\small\bfseries\color{#1}#2\par}\vspace{3pt}
    {\small\raggedright
      \setlength{\abovedisplayskip}{4pt}
      \setlength{\belowdisplayskip}{4pt}
      \setlength{\abovedisplayshortskip}{3pt}
      \setlength{\belowdisplayshortskip}{3pt}
      #3\par}
    \end{minipage}
  \endgroup\par\vspace{9pt}}
% Figura (si existe) a la izquierda y condiciones del resultado al costado [GPRES, p. 1].
\newcommand{\grafico}[3][6.4cm]{%
  \IfFileExists{figuras/#2}{\includegraphics[width=\linewidth,height=#1,keepaspectratio]{#2}}%
  {\fcolorbox{accent}{white}{\parbox[c][5.4cm][c]{.92\linewidth}{\centering\color{accent}Figura pendiente:\\\texttt{#2}\\#3}}}}
\newcommand{\figura}[3][.68]{%
  \begin{columns}[c,onlytextwidth]
    \begin{column}{#1\textwidth}\centering\grafico{#2}{}\end{column}
    \begin{column}{\dimexpr.98\textwidth-#1\textwidth\relax}\small #3\end{column}
  \end{columns}}
% Fotograma y enlace en el PDF; en el PPTX se reemplaza por el video [GPRES, p. 3].
\newcommand{\animacion}[2]{%
  \begin{columns}[c,onlytextwidth]
    \begin{column}{.74\textwidth}\centering
      \grafico{anim-#1.png}{}\par\smallskip
      {\small\color{accent}\enlace{#1}}
    \end{column}
    \begin{column}{.24\textwidth}\small #2\end{column}
  \end{columns}}

\tikzset{
  flecha/.style={-{Stealth[length=5pt]},thick,itbablue},
  caja/.style={draw=itbablue,thick,rounded corners=3pt,fill=softblue,align=center,font=\small,inner sep=5pt},
}

\title[TP5: evacuación híbrida]{Dinámica peatonal: evacuación con esquema híbrido}
\subtitle{Event Driven + Time Step Driven con el Social Force Model}
\author[Grupo 07]{Grupo 07 - Comisión S\\Eugenio Migliaro, Francisco Costa y Franco Branda}
\institute{72.25 Simulación de Sistemas - ITBA}
\date{16 de noviembre de 2026}

\begin{document}

\begin{frame}\titlepage\end{frame}

%=====================================================================
\section{Introducción}
\begin{frame}\sectionpage\end{frame}

\begin{frame}{Evacuación por una puerta angosta}
\begin{columns}[c,onlytextwidth]
\begin{column}{.47\textwidth}\centering
\begin{tikzpicture}[scale=.082,>=Stealth]
  \draw[very thick] (0,0) -- (24.6,0); \draw[very thick] (25.4,0) -- (50,0);
  \draw[very thick] (0,0) -- (0,50) -- (50,50) -- (50,0);
  \foreach \x/\y in {8/40,15/33,30/42,41/30,22/24,36/15,10/14,44/45,27/36,5/27,47/8,18/6} {\fill[evento] (\x,\y) circle (0.8);}
  \foreach \x/\y in {23.2/1.2,25/1.4,26.8/1.2,24.1/2.8,25.9/2.8,22.4/3,27.6/3,25/4.3} {\fill[paso] (\x,\y) circle (0.8);}
  \draw[<->,thin] (0,53) -- (50,53) node[midway,above] {$L=50$ m};
  \draw[->,accent] (34,-3) node[right,font=\scriptsize] {puerta, $d=0{,}8$ m} -- (25.6,-0.5);
  \draw[dashed,gray] (19,-8) -- (31,-8) node[right,font=\scriptsize,black] {se retira a 1 m};
\end{tikzpicture}
\end{column}
\begin{column}{.51\textwidth}
\ficha{Sistema}{%
$N=500$ agentes en una sala cuadrada, con una única puerta en el centro de una pared. Arrancan al azar, sin superponerse, y quieren salir a $v_d=1{,}5$ m/s.}
\ficha{Dos regímenes en la misma evacuación}{%
\textcolor{evento}{Diluido}: van en línea recta y casi no interactúan.\\
\textcolor{paso}{Cuello de botella}: contactos continuos frente a la puerta.}
\ficha{Pregunta}{%
¿Se puede resolver cada agente con el esquema que le corresponde, por eventos o por paso temporal, sin cambiar la física?}
\end{column}
\end{columns}
\end{frame}

\begin{frame}{Social Force Model sin fuerza social}
\begin{columns}[T,onlytextwidth]
\begin{column}{.55\textwidth}
\ficha{Ecuación de movimiento [Helbing et al. 2000]}{%
\[m_i\frac{d\vect v_i}{dt}=m_i\frac{v_d\,\vect e_i-\vect v_i}{\tau}+\sum_{j}\vect f_{ij}+\sum_{W}\vect f_{iW}\]}
\ficha{Contacto: solo si $\xi_{ij}=r_i+r_j-d_{ij}>0$}{%
\[\vect f_{ij}=k\,\xi_{ij}\,\vect n_{ij}+\kappa\,\xi_{ij}\,\big[(\vect v_j-\vect v_i)\cdot\vect t_{ij}\big]\,\vect t_{ij}\]
$\vect n_{ij}$: de $j$ hacia $i$; $\vect t_{ij}$: perpendicular. Paredes análogas; los marcos de la puerta, como puntos.}
\ficha{Sin fuerza social}{%
Nadie reacciona a otro hasta tocarlo. Un agente libre y relajado va en línea recta a $v_d$.}
\end{column}
\begin{column}{.42\textwidth}
\ficha{Parámetros}{%
\renewcommand{\arraystretch}{1.15}
\begin{tabular}{@{}l@{\hspace{10pt}}l@{}}
$\tau$ & 0,5 s\\
$k$ & $1{,}2\times10^{5}$ N/m\\
$\kappa$ & $2{,}4\times10^{5}$ kg/(m\,s)\\
$r_i$ & uniforme en $[0{,}18;\,0{,}25]$ m\\
$m_i$ & $80\text{ kg}\,(r_i/0{,}215\text{ m})^2$
\end{tabular}}
\ficha[accent]{Masa proporcional al área}{%
El paper usa 80 kg para todos: es un apartamiento propio.}
\ficha{Targets}{%
Adentro: el punto del 60\,\% central de la puerta más cercano. Al cruzarla: alejarse de la puerta ($-\hat{\vect y}$).}
\end{column}
\end{columns}
\end{frame}

%=====================================================================
\section{Implementación}
\begin{frame}\sectionpage\end{frame}

\begin{frame}{TSDMD: Beeman predictor-corrector}
\begin{columns}[T,onlytextwidth]
\begin{column}{.56\textwidth}
\ficha{Fuerzas que dependen de la velocidad}{%
\[\vect r_{n+1}=\vect r_n+\dt\,\vect v_n+\dt^2\big(\tfrac23\vect a_n-\tfrac16\vect a_{n-1}\big)\]
\[\vect v^{\,p}_{n+1}=\vect v_n+\dt\big(\tfrac32\vect a_n-\tfrac12\vect a_{n-1}\big)\]
\[\vect v_{n+1}=\vect v_n+\dt\big(\tfrac13\vect a(\vect r_{n+1},\vect v^{\,p}_{n+1})+\tfrac56\vect a_n-\tfrac16\vect a_{n-1}\big)\]
$\vect a_{n+1}$ se reevalúa con $\vect v_{n+1}$, reusando los contactos del paso.}
\end{column}
\begin{column}{.40\textwidth}
\ficha{Contactos}{Celdas de lado $\geq2r_{\max}$ (CIM): cada par de celdas vecinas una sola vez.}
\ficha{Paso temporal}{$\dt=\valor{dtElegido}$ ms, elegido con un test de convergencia (choque aislado y caudal).}
\ficha{Salida}{Cuadros cada $\dt_2$ y el instante de salida de cada agente, interpolado dentro del paso.}
\end{column}
\end{columns}
\end{frame}

\begin{frame}{EDMD puro: discos duros entre eventos}
\begin{columns}[T,onlytextwidth]
\begin{column}{.50\textwidth}
\ficha{Eventos}{%
\begin{itemize}\setlength\itemsep{1pt}
\item choque entre agentes (cuadrática del TP3)
\item choque con pared o marco de la puerta
\item cruce de celda (busca compañeros nuevos)
\item cruce de la puerta: segundo target
\item retiro a 1 m de la puerta
\item patada de relajación, cada $\delta$
\end{itemize}}
\ficha{Cola de prioridad}{Invalidación perezosa con un contador por agente, como en el TP3.}
\end{column}
\begin{column}{.46\textwidth}
\ficha{Choque instantáneo derivado del SFM}{%
El resorte normal no amortigua: la normal es \textbf{elástica}. La fricción, integrada en medio período del resorte:
\[\frac{v_t'}{v_t}=e_t=\exp\!\Big(-\frac{2\kappa\,|v_n|}{k}\Big)\]
La masa reducida se cancela: vale para masas distintas y contra paredes.}
\ficha[accent]{Recuperar $v_d$ después del choque}{%
Patadas: $\vect v\leftarrow v_d\vect e+(\vect v-v_d\vect e)\,e^{-\delta/\tau}$, recta entre eventos.}
\end{column}
\end{columns}
\end{frame}

\begin{frame}{Recuperar $v_d$ en EDMD: patadas o relajación analítica}
\begin{columns}[T,onlytextwidth]
\begin{column}{.48\textwidth}
\ficha{Patadas cada $\delta$}{%
Recta entre eventos: los choques siguen siendo una cuadrática. A todos los agentes sin relajar, cada $\delta$:
\[\vect v\leftarrow v_d\vect e+(\vect v-v_d\vect e)\,e^{-\delta/\tau}\]
Es la relajación exacta del SFM durante $\delta$, con $\vect e$ fija.}
\end{column}
\begin{column}{.48\textwidth}
\ficha{Relajación analítica}{%
\[\vect x(t)=\vect x_0+\vect u\,s+\tau\,\vect w\,\big(1-e^{-s/\tau}\big),\qquad\vect u=v_d\vect e\]
Eventos propios: fin de la relajación ($|\vect w|e^{-s/\tau}\leq0{,}01\,v_d$) y reapuntado, cuando la dirección al target se aparta de la congelada más que una tolerancia.}
\ficha[accent]{Costo}{Los choques ya no son una cuadrática: avance conservador de segundo orden, $g(t+h)\geq g+\dot g\,h-A h^2/2$.}
\end{column}
\end{columns}
\end{frame}

\begin{frame}{Híbrido (1): un esquema por agente}
\begin{columns}[c,onlytextwidth]
\begin{column}{.42\textwidth}\centering
\begin{tikzpicture}
  \node[caja,draw=evento,fill=evento!10,minimum width=3.2cm] (e) {\textcolor{evento}{\textbf{EDMD}}\\recto, $\vect v=v_d\vect e$};
  \node[caja,draw=paso,fill=paso!10,minimum width=3.2cm,below=2.4cm of e] (p) {\textcolor{paso}{\textbf{TSDMD}}\\SFM con Beeman};
  \draw[flecha,paso] ([xshift=8mm]e.south west) -- node[left,font=\scriptsize,align=right,paso] {contacto\\posible en\\menos de $\tstar$} ([xshift=8mm]p.north west);
  \draw[flecha,evento] ([xshift=-8mm]p.north east) -- node[right,font=\scriptsize,align=left,evento] {aislado y\\relajado, sin\\contacto posible\\en $2\tstar$} ([xshift=-8mm]e.south east);
\end{tikzpicture}
\end{column}
\begin{column}{.56\textwidth}
\ficha{Contacto posible}{%
Contra EDMD y paredes: exacto (cuadrática).\\
Contra TSDMD, cuyo futuro no se conoce, con la cota $|\vect v_j|\leq V_{\max}=2v_d$:
\[t_{\rm seg}=\frac{d_{ij}-(r_i+r_j)}{|\vect v_i|+V_{\max}}\]}
\ficha{Parámetro principal: $\tstar$}{%
$\tstar\to\infty$: TSDMD puro. $\tstar$ chico: cada agente queda en EDMD hasta el último momento. Requiere $\tstar\geq\dt$.}
\ficha{Histéresis}{Volver exige $2\tstar$: evita pases de ida y vuelta.}
\end{column}
\end{columns}
\end{frame}

\begin{frame}{Híbrido (2): la grilla de $\dt$ es el reloj maestro}
\centering
\begin{tikzpicture}[x=1.25cm,>=Stealth]
  \draw[thick,gray] (0,0) -- (10.4,0);
  \foreach \i/\l in {0/{t_0},2/{t_1},4/{\cdots},6/{t_n},8/{t_{n+1}},10/{\cdots}} {\draw[paso,very thick] (\i,-.25) -- (\i,.25); \node[paso,below=4pt] at (\i,-.2) {$\l$};}
  \foreach \x in {.7,1.3,3.1,6.6,7.2,7.7,9.1} {\fill[evento] (\x,.45) -- ++(-.12,.2) -- ++(.24,0) -- cycle;}
  \draw[|-|,thick,itbablue] (6,1.0) -- (8,1.0) node[midway,above=2pt,font=\small] {un intervalo};
\end{tikzpicture}
\vspace{6pt}
\begin{columns}[T,onlytextwidth]
\begin{column}{.52\textwidth}
\ficha{En cada intervalo $(t_n,t_{n+1}]$}{%
\begin{enumerate}\setlength\itemsep{1pt}
\item eventos EDMD en orden: alertas $\tstar$ antes de un contacto posible, cruce de puerta, retiro
\item un paso de Beeman para los TSDMD
\item en $t_{n+1}$: revisiones de la cota y cambios de esquema
\end{enumerate}}
\end{column}
\begin{column}{.44\textwidth}
\ficha{Agentes EDMD}{No se integran: $\vect x(t)=\vect x(t_0)+\vect v\,(t-t_0)$ se evalúa cuando hace falta (salida, criterio).}
\ficha{TSDMD}{Solo interactúan entre ellos y con las paredes: nunca hay un EDMD en contacto.}
\end{column}
\end{columns}
\end{frame}

\begin{frame}{Híbrido (3): no perder contactos en el pase}
\begin{columns}[T,onlytextwidth]
\begin{column}{.50\textwidth}
\ficha{Revisión de la cota}{%
Siempre en instantes de la grilla, con las posiciones TSDMD exactas, y al menos un paso después de la anterior: si $t_{\rm seg}\leq\tstar$ pasa a TSDMD ya; si no, se revisa de nuevo en
\[t_n+\Big\lfloor\frac{t_{\rm seg}-\tstar}{\dt}\Big\rfloor\dt\]}
\ficha{Garantía}{Con $\tstar\geq\dt$ el pase ocurre antes de cualquier contacto posible: el contacto siempre lo resuelve el SFM. $V_{\max}$ se verifica en cada paso: en (a) la rapidez máxima fue \valor{rapidezMaxima}; violaciones, \valor{violacionesCota}.}
\end{column}
\begin{column}{.46\textwidth}
\ficha{EDMD $\to$ TSDMD en $t_n$}{%
$\vect x(t_n)$ de la recta; $\vect a_n=\vect a_{n-1}=0$, exacto porque iba relajado y sin contactos.}
\ficha{TSDMD $\to$ EDMD en $t_n$}{%
Relajado si $|\vect v-v_d\vect e|\leq0{,}01\,v_d$; se redondea $\vect v=v_d\vect e$ y se predicen sus eventos.}
\ficha[accent]{Por etapas}{Al principio casi todos van por eventos; en el tapón, todos por paso temporal.}
\end{column}
\end{columns}
\end{frame}

\begin{frame}{Híbrido (4): choques y recuperación de $v_d$}
\begin{columns}[T,onlytextwidth]
\begin{column}{.50\textwidth}
\ficha{Choques}{%
En el híbrido no hay choques instantáneos: un agente EDMD pasa a TSDMD $\tstar$ antes de cualquier contacto posible con otro agente o con una pared. Todos los contactos son blandos, del SFM.}
\ficha{Recuperación de $v_d$}{%
La hace el término de deseo del SFM, integrado con Beeman. El agente vuelve a EDMD cuando está relajado y aislado.}
\end{column}
\begin{column}{.46\textwidth}
\ficha{¿Conviene relajar en forma analítica?}{%
Solo ahorraría los pasos de agentes aislados pero sin relajar: \valor{aisladosSinRelajar} de los pasos TSDMD. No se implementó.}
\ficha{Cota del ahorro}{%
El híbrido solo ahorra el vuelo libre: \valor{fraccionTSDMD} del tiempo-agente es TSDMD, porque el tapón dura casi toda la evacuación.}
\end{column}
\end{columns}
\end{frame}

\begin{frame}{Arquitectura del motor (C++)}
\begin{columns}[T,onlytextwidth]
\begin{column}{.52\textwidth}
\ficha{Módulos}{%
\renewcommand{\arraystretch}{1.1}
\begin{tabular}{@{}l@{\hspace{8pt}}p{.62\linewidth}@{}}
\texttt{model} & sala, targets y parámetros\\
\texttt{setup} & estado inicial al azar\\
\texttt{sfm} & fuerzas con celdas y Beeman\\
\texttt{events} & predicción y regla de choque\\
\texttt{tsdmd} & paso temporal puro\\
\texttt{edmd} & eventos puro\\
\texttt{hybrid} & esquema híbrido\\
\texttt{output} & salida de texto
\end{tabular}}
\end{column}
\begin{column}{.44\textwidth}
\ficha{Salida: solo estado}{%
Cuadros \texttt{x y vx vy modo} cada $\dt_2$, instantes de salida y cambios de esquema. Los observables se calculan aparte, en Python.}
\ficha{Validación}{El híbrido con $\tstar$ enorme da exactamente las mismas salidas que el TSDMD puro.}
\end{column}
\end{columns}
\end{frame}

%=====================================================================
\section{Simulaciones}
\begin{frame}\sectionpage\end{frame}

\begin{frame}{Estudios realizados}
\centering\small
\renewcommand{\arraystretch}{1.45}
\setlength{\tabcolsep}{4pt}
\begin{tabular}{@{}llll@{}}
\toprule
\textcolor{itbablue}{\textbf{Estudio}} & Esquema & Variable & \textcolor{itbablue}{\textbf{Realizaciones}}\\
\midrule
Paso temporal & TSDMD & $\dt\in\{0{,}5;\,1;\,2\}$ ms & 10 por $\dt$\\
Atascos & TSDMD & fluctuación en contacto & 20 y 10 por amplitud\\
Segundo target & TSDMD & punto a $D$ o alejarse & 10 por opción\\
(a) y (b) & los tres & esquema & 20 por esquema\\
(c) & híbrido & $\tstar$ de 0,002 a 1 s & 10 por $\tstar$\\
(d) & los tres & $N$ de 100 a 1000, dos variantes & 5 por $N$ (EDMD: 3)\\
Recuperación de $v_d$ & EDMD & $\delta$ y tolerancia de reapuntado & 10 ($N=60$, $L=15$ m)\\
\bottomrule
\end{tabular}
\par\medskip
\footnotesize Mismas condiciones iniciales en todas las variantes de cada estudio. Corte: 300 s sin salidas (atasco).
\end{frame}

\begin{frame}{Observables}
\small
\begin{columns}[T,onlytextwidth]
\begin{column}{.49\textwidth}
\ficha{Curva de población y evacuación}{%
$N_{\rm adentro}(t)=N-\#\{\text{salidas}\leq t\}$\\
$t_{\rm evac}$: salida del último agente, si la corrida termina.}
\ficha{Caudal}{$J$: salidas entre el 10 y el 90\,\% sobre el tiempo que tardan.}
\ficha{Atasco}{Una corrida se traba si no sale nadie durante 300 s. Se informa la fracción, con intervalo de Clopper--Pearson del 95\,\%.}
\end{column}
\begin{column}{.49\textwidth}
\ficha{Esquemas}{Fracción de los agentes adentro resueltos por EDMD y por TSDMD, en cada instante.}
\ficha{Tiempo de ejecución}{Wall-clock del bucle, de las corridas completas. Complemento: costo por segundo simulado, con todas las corridas.}
\ficha{Promedios entre las $M$ realizaciones}{Cada realización aporta un valor; media y desvío estándar entre realizaciones. $t_{\rm evac}$: solo de las que terminan.}
\end{column}
\end{columns}
\end{frame}

%=====================================================================
\section{Resultados}
\begin{frame}\sectionpage\end{frame}

\begin{frame}{Validación de la regla de choque}
\figura{regla-choque.pdf}{
\ficha{Choque aislado}{Dos agentes sin término de deseo, con el SFM blando y $\dt=0{,}1$ ms, contra la regla de disco duro.}
\ficha{Resultado}{La tangencial coincide en todo el rango; en choques rasantes ($b=0{,}8\,\sigma$) la normal pierde hasta 5\,\%.}}
\end{frame}

\begin{frame}{Paso temporal}
\begin{columns}[c,onlytextwidth]
\begin{column}{.48\textwidth}\centering\grafico{dt-choque.pdf}{}\end{column}
\begin{column}{.48\textwidth}\centering\grafico{dt-caudal.pdf}{}\end{column}
\end{columns}
\par\small Izquierda: choque frontal de los dos agentes más livianos (56 kg). Derecha: $N=500$, 10 realizaciones por $\dt$; el caudal no cambia: $\dt=\valor{dtElegido}$ ms.
\end{frame}

\begin{frame}{Atascos permanentes}
\animacion{atasco}{%
\ficha{TSDMD, $N=500$}{Arco estático en la puerta: dos agentes apoyados en la pared sostienen al resto.}
\ficha{Por qué no se rompe}{La fricción del SFM es viscosa: en reposo no actúa, y no hay fluctuaciones.}}
\end{frame}

\begin{frame}{Atascos: por qué a veces sí y a veces no}
\figura{compresion.pdf}{
\ficha{TSDMD}{Contactos a menos de 1,5 m de la puerta.}
\ficha[accent]{Hipótesis: compresión}{Con la sala llena, los contactos se comprimen hasta 6 cm y los agentes se cuelan. Con pocos adentro, 1 cm: el arco aguanta.}
\ficha{Prueba}{Con $k$ diez veces mayor, 3 semillas que evacuaban se traban al principio.}}
\end{frame}

\begin{frame}{Atascos: curvas de población}
\figura{curvas-poblacion-20-semillas.pdf}{
\ficha{TSDMD, 20 realizaciones}{Se traban \valor{atascosSinFluctuacion}. Las que evacúan tardan $\valor{tevacSinFluctuacion}$ s.}
\ficha{Corte}{300 s sin salidas: la pausa más larga que terminó rompiéndose fue de 121 s.}}
\end{frame}

\begin{frame}{Atascos y fluctuaciones}
\figura{probabilidad-atasco.pdf}{
\ficha{Fuerza aleatoria en contacto}{Gaussiana, de amplitud fija como fracción de $m v_d/\tau$, renovada cada 0,1 s, solo sobre agentes que se tocan.}
\ficha{Resultado}{No reduce los atascos de forma detectable: se mantiene el modelo sin fluctuación.}}
\end{frame}

\begin{frame}{Segundo target}
\figura{segundo-target.pdf}{
\ficha{TSDMD, 10 realizaciones por opción}{Target común en $(L/2,-D)$ contra alejarse de la puerta.}
\ficha{Resultado}{Todas dentro de las barras: el cuello de botella está adentro. Se usa alejarse.}}
\end{frame}

\begin{frame}{(a) TSDMD puro}
\animacion{tsdmd}{\ficha{$N=500$, $\dt=\valor{dtElegido}$ ms}{Todos los agentes, en todos los pasos.}}
\end{frame}

\begin{frame}{(a) EDMD puro}
\animacion{edmd}{\ficha{Discos duros, patadas cada $\delta=0{,}05$ s}{Se traba en las 20 corridas, al principio: salen \valor{salidasEDMD} agentes.}
\ficha{Costo (punto c)}{\valor{costoEDMD} ms de reloj por segundo simulado, contra \valor{costoTSDMD} con TSDMD: \valor{costoRelativoEDMD} veces más.}}
\end{frame}

\begin{frame}{(a) Híbrido}
\animacion{hibrido}{\ficha{$\tstar=\valor{tstarElegido}$ s}{\textcolor{evento}{EDMD} en vuelo libre; \textcolor{paso}{TSDMD} cerca de un contacto posible y en el tapón.}}
\end{frame}

\begin{frame}{(a) Fracción de agentes en cada esquema}
\figura{fraccion-esquemas.pdf}{
\ficha{Híbrido, $\tstar=\valor{tstarElegido}$ s}{Fracción de los agentes adentro, media de \valor{realizacionesPoblacion} realizaciones.}
\ficha{Resultado}{Todos pasan a TSDMD en \valor{tiempoVueloLibre}: después, el tapón.}}
\end{frame}

\begin{frame}{(b) Curvas de población}
\figura{poblacion-esquemas.pdf}{
\ficha{$N=500$, \valor{realizacionesPoblacion} realizaciones por esquema}{Media y desvío entre realizaciones de las que terminan; las trabadas, aparte.}}
\end{frame}

\begin{frame}{(b) Tiempo de evacuación y atascos}
\figura[.72]{tevac-esquemas.pdf}{
\ficha{De las que terminan}{\valor{textoEvacuacion}}
\ficha{Atascos}{Intervalo de Clopper--Pearson del 95\,\%.}}
\end{frame}

\begin{frame}{(c) Tiempo de ejecución en función de $\tstar$}
\centering\grafico[5.5cm]{tiempos-tstar.pdf}{}\par\smallskip
\small $N=500$, 10 realizaciones por $\tstar$ (terminan \valor{terminanc}), en la misma computadora: \valor{maquina}. Líneas: TSDMD y EDMD puros (media y desvío).
\end{frame}

\begin{frame}{(c) Tiempo de evacuación en función de $\tstar$}
\figura{tevac-tstar.pdf}{
\ficha{Física}{Si el pase no pierde contactos, $t_{\rm evac}$ no depende de $\tstar$.}
\ficha{Atascos}{Abajo, cuántas corridas terminan con cada $\tstar$.}}
\end{frame}

\begin{frame}{(c) Comparación con TSDMD semilla por semilla}
\begin{columns}[T,onlytextwidth]
\begin{column}{.48\textwidth}
\ficha{$\tstar\geq0{,}5$ s: nadie es EDMD}{Mismas salidas que TSDMD en \valor{mismasSalidas} semillas. El híbrido cuesta \valor{sobrecostoHibrido} más: es el costo fijo del esquema (eventos, revisiones de la cota).}
\ficha{$\tstar\leq0{,}2$ s}{Costo por segundo simulado sobre el de TSDMD de la misma semilla: de \valor{rangoPareado} en promedio, con desvíos del 10 al 35\,\%. Las trayectorias divergen y ninguna ganancia se distingue.}
\end{column}
\begin{column}{.48\textwidth}
\ficha[accent]{Por qué}{El ahorro posible es el vuelo libre: a lo sumo el 2,6\,\% del tiempo-agente (\valor{fraccionTSDMD} es TSDMD). Es del orden del costo fijo y mucho menor que la dispersión entre realizaciones.}
\ficha{EDMD puro}{\valor{costoRelativoEDMD} veces el costo de TSDMD.}
\end{column}
\end{columns}
\end{frame}

\begin{frame}{(d) Tiempo de ejecución en función de $N$}
\centering\grafico[5.8cm]{escalamiento-wall.pdf}{}\par\smallskip
\small Corridas que terminan, de 5 realizaciones por $N$ (EDMD puro solo termina con $N=100$). Izquierda: densidad inicial constante, $L=50\text{ m}\sqrt{N/500}$. Derecha: $L=50$ m.
\end{frame}

\begin{frame}{(d) Costo por segundo simulado en función de $N$}
\centering\grafico[5.5cm]{escalamiento-costo.pdf}{}\par\smallskip
\small Todas las corridas, incluidas las trabadas; EDMD puro: 3 realizaciones por $N$, hasta $N=600$. Ajuste $N^{\alpha}$ con densidad constante: TSDMD e híbrido $\alpha=\valor{exponenteDensidadTSDMD}$, EDMD $\alpha=\valor{exponenteDensidadEDMD}$; con $L=50$ m: $\valor{exponenteSalaTSDMD}$ y $\valor{exponenteSalaEDMD}$.
\end{frame}

\begin{frame}{Recuperación de $v_d$ en EDMD puro}
\figura[.6]{recuperacion.pdf}{
\ficha{Sistema chico}{$N=60$, $L=15$ m, 10 realizaciones; diferencia de caudal con el TSDMD de la misma semilla.}
\ficha{Resultado}{\valor{textoRecuperacion}}
\ficha[accent]{Conviene}{En EDMD puro, patadas. En el híbrido la recupera el SFM, al costo de TSDMD.}}
\end{frame}

%=====================================================================
\section{Conclusiones}
\begin{frame}\sectionpage\end{frame}

\begin{frame}{Conclusiones}
\begin{columns}[T,onlytextwidth]
\begin{column}{.48\textwidth}
\ficha{Modelo}{Con la puerta de 0,8 m, el SFM sin fuerza social se traba en cerca del 40\,\% de las corridas: arcos estáticos que la fricción viscosa no rompe. La compresión de los contactos explica por qué casi siempre se traban al final.}
\ficha{Esquema híbrido}{Reproduce el TSDMD: mismo tiempo de evacuación, mismo caudal y la misma fracción de atascos, sin violaciones de la cota. Los agentes van por eventos solo en los primeros 30 s.}
\end{column}
\begin{column}{.48\textwidth}
\ficha{EDMD puro}{Con discos duros se traba siempre, al principio, y cuesta \valor{costoRelativoEDMD} veces más que TSDMD: los contactos permanentes se vuelven millones de choques. Para recuperar $v_d$, las patadas son más baratas y más fieles que la relajación analítica.}
\ficha{Costo}{El ahorro del híbrido está acotado por el tapón (\valor{fraccionTSDMD} del tiempo-agente es TSDMD) y no se distingue de su costo fijo: en la práctica cuesta lo mismo que TSDMD. EDMD puro cuesta \valor{costoRelativoEDMD} veces más y escala peor con $N$.}
\end{column}
\end{columns}
\end{frame}

\begin{frame}{Declaración de uso de IA}
\begin{columns}[T,onlytextwidth]
\begin{column}{.40\textwidth}
\ficha{Herramienta}{Claude Code con Claude Opus 5.5.}
\ficha{Intervención}{Código C++ y Python, pruebas, corridas, gráficos, animaciones y borrador de la presentación. Las decisiones de diseño las tomó el grupo.}
\end{column}
\begin{column}{.56\textwidth}
\ficha[accent]{Decisiones en las que el grupo modificó la recomendación del agente}{%
\footnotesize
\begin{itemize}\setlength\itemsep{0pt}
\item Modelo: SFM en lugar de CPM.
\item Masa proporcional al área en lugar de 80 kg para todos.
\item Estudio formal del segundo target, con más posiciones, usando la mejor opción en lugar de fijarla de antemano.
\item Comparar patadas y relajación analítica en lugar de usar solo patadas.
\item Presentación sin límite estricto de 13 minutos.
\end{itemize}}
\end{column}
\end{columns}
\end{frame}

\end{document}
